\documentclass[10pt,a4paper]{amsart}
\usepackage[margin=19mm]{geometry}
\usepackage{amsmath,amssymb,mathtools}
\usepackage[scr=rsfs,cal=euler]{mathalfa}
\usepackage{xfrac}
\usepackage[unicode,hidelinks,bookmarks=false]{hyperref}
\newcommand{\ie}{i.e.\ }
\newcommand{\cf}{cf.\ }
\newcommand{\resp}{respectively\ }
\newcommand{\Subsection}{\subsection}
\providecommand{\nobreakdash}{\nobreak\hbox{-}}
\setlength{\emergencystretch}{2em}
\multlinegap0pt
\usepackage{bm}
\usepackage{bbm}
\usepackage{dsfont}
\usepackage{xspace}
\usepackage{cancel}
\usepackage{mathtools}
\usepackage{amsthm}
\makeatletter
\@ifundefined{theorem}{\newtheorem{theorem}{Theorem}}{}
\makeatother
\DeclareMathOperator{\val}{val}
\title[Computing Expansions in Infinitely Many Cantor Real Bases vi]{Computing Expansions in Infinitely Many Cantor Real Bases via a Single Transducer}
\author{\'Emilie Charlier, Pierre Popoli, Michel Rigo}
\date{Source: arXiv:2507.04848v1 (CC BY 4.0)}
\begin{document}
\maketitle

\noindent\emph{Source and licence.} Computing Expansions in Infinitely Many Cantor Real Bases via a Single Transducer. Version arXiv:2507.04848v1 (CC BY 4.0), distributed under the Creative Commons Attribution 4.0 International licence, as recorded in the arXiv licence metadata for this exact version and shown on the abstract page https://arxiv.org/abs/2507.04848v1. Full rights notice: licenses/arxiv-2507.04848-CC-BY-4.0.txt.\par\smallskip

\noindent\textit{Editorial scope.} Counted unit: the Theorem whose source label is \texttt{thm:TM23} in the article (source0.tex lines 190--196, source label \texttt{thm:TM23}), delivered together with the article's own complete original proof of it (source0.tex lines 198--250, one \texttt{proof} environment of 53 lines). No mathematics is written, repaired or re-derived: statement and proof are the article's own text line for line. Required context retained uncounted: none. Every definition, display and label of those lines is retained unchanged. Omitted from this package and not redistributed: 1--189, 197--197, 251--1835. Nothing inside a retained excerpt is omitted or altered: every retained body is delivered exactly as the article's lines are written, so no omission receipt of any kind accompanies this package. Citations: the retained excerpts carry no citation command at all, so no bibliography entry of the article is transcribed and no reference is redistributed. Reference adaptation: none was needed and none was made; every reference inside the delivered unit resolves inside it, so no unresolved reference or citation is produced. Printed slips kept literal: no printed word, symbol, hypothesis or number of the counted statement or of its proof was altered. Layout and class: the article's own class is \texttt{amsart} with packages including geometry, amsaddr, amsmath, amssymb, amsthm, algorithm, algorithmicx, algpseudocode, enumitem, thmtools, hyperref, xcolor, graphicx, cleveref; the wrapper is the lane's pinned \texttt{amsart} editorial wrapper at 10pt on A4, which restates the theorem-like environments the retained text needs and adds the article's own macro definitions that the retained text needs (\texttt{\textbackslash val}), with extra wrapper packages (bm, bbm, dsfont, xspace, cancel, mathtools). The delivered numbering is the unit's own. Assets: no retained span contains a figure environment, a graphics-inclusion command, a PDF-page inclusion, a table or a TikZ picture; no image file is copied into this package, the package declares no asset dependency and no image file is redistributed with it. Licence: version arXiv:2507.04848v1 of this article is distributed under the Creative Commons Attribution 4.0 International licence, as recorded in the arXiv licence metadata for this exact version and shown on the abstract page https://arxiv.org/abs/2507.04848v1. A full rights notice is delivered at licenses/arxiv-2507.04848-CC-BY-4.0.txt. Characters: every retained excerpt is pure ASCII, so the delivered engine, XeLaTeX with its default Latin Modern text font, reports no missing character.
\begin{theorem} \label{thm:TM23}
For all $r\in[0,1)$, the ${\mathbf{T}}$-expansion of $r$ is the sequence 
\[
    h_2(c_0)h_3(c_1)h_3(c_2)h_2(c_3)\cdots 
\]
obtained from $d_6(r)=c_0c_1c_2c_3\cdots$ by application of the morphisms $h_2$ and $h_3$ in the ordering given by the Thue--Morse word~${\mathbf{T}}$. 
\end{theorem}

\begin{proof}
Write ${\mathbf{T}}=(\beta_n)_{n\in\mathbb{N}}$, and let 
\[
    \mathbf{y}=y_0y_1y_2\cdots=h_2(c_0)h_3(c_1)h_3(c_2)h_2(c_3)\cdots h_{\beta_n}(c_n)\cdots.
\]
Otherwise stated, we set $h_{\beta_n}(c_n)=y_{2n}y_{2n+1}$ for all $n\ge 0$. We show that $\val_{\mathbf{T}}(\mathbf{y})=r$ and 
\[
    \forall\ell\ge 1,\quad  \sum_{n=\ell}^{+\infty} \frac{y_n}{\prod_{i=0}^n\beta_i}
    <\frac{1}{\prod_{i=0}^{\ell-1} \beta_i}
\]
proving that $\mathbf{y}$ is indeed get the ${\mathbf{T}}$-expansion of $r$. Let us first prove that $\mathbf{y}$ is a $\mathbf{T}$-representation of $r$. We have
\begin{eqnarray*}
    \val_{\mathbf{T}}(\mathbf{y}) 
    &=& \sum_{n= 0}^{+\infty}\left( \frac{y_{2n}}{\prod_{i=0}^{2n}\beta_i}
    +\frac{y_{2n+1}}{\prod_{i=0}^{2n+1}\beta_i}\right)
    =\sum_{n=0}^{+\infty}\frac{1}{6^{n+1}} (y_{2n}\beta_{2n+1}+y_{2n+1}).
\end{eqnarray*}

It suffices to show that $c_n=y_{2n}\beta_{2n+1}+y_{2n+1}$.
We have two cases to consider regarding the value of $\beta_n$. If $\beta_n=2$ then, on the one hand, $c_n= y_{2n}\cdot 3+y_{2n+1}$ and on the other hand, $\tau(\beta_n)=\beta_{2n}\beta_{2n+1}=23$, hence $\beta_{2n+1}=3$. Similarly, if $\beta_n=3$ then $c_n= y_{2n}\cdot 2+y_{2n+1}$ and $\beta_{2n+1}=2$ since $\tau(\beta_n)=\beta_{2n}\beta_{2n+1}=32$.

We now turn to the greediness of the $\mathbf{T}$-representation. Let us first assume that the tail is starting with an even index~$2\ell$. Then
\[
    \sum_{n= 2\ell}^{+\infty} \frac{y_n}{\prod_{i=0}^n\beta_i}
    =\sum_{n= \ell}^{+\infty} \left( \frac{y_{2n}}{\prod_{i=0}^{2n}\beta_i}
    +\frac{y_{2n+1}}{\prod_{i=0}^{2n+1}\beta_i}\right)
    =\sum_{n= \ell}^{+\infty} \frac{c_n}{6^{n+1}}
    <\frac{1}{6^{\ell}}
    =\frac{1}{\prod_{i=0}^{2\ell-1} \beta_i},
\]
where we used the fact that $(c_n)_{n\in\mathbb{N}}$ is the greedy $\delta$-expansion of $r$. Assume now that the tail is starting with an odd index~$2\ell+1$. We get
\[
    \sum_{n=2\ell+1}^{+\infty} \frac{y_n}{\prod_{i=0}^n\beta_i}
    =\frac{y_{2\ell+1}}{\prod_{i=0}^{2\ell+1}\beta_i} 
    + \sum_{n=\ell+1}^{+\infty} \left( \frac{y_{2n}}{\prod_{i=0}^{2n}\beta_i}
    +\frac{y_{2n+1}}{\prod_{i=0}^{2n+1}\beta_i}\right)
    =\frac{y_{2\ell+1}}{6^\ell} +\sum_{n\ge \ell+1} \frac{c_n}{6^{n+1}}.
\]
Again, we have two cases. If $\beta_\ell=2$ then 
$y_{2\ell+1}\le 2$ by definition of the morphism $h_2$. Using that $(c_n)_{n\in\mathbb{N}}$ is a $\delta$-expansion and that $\beta_{2\ell}=2$, we see that the above quantity is less than
\[
    \frac{2}{6^{\ell+1}}+\frac{1}{6^{\ell+1}}
    =\frac{3}{2\cdot 3\cdot 6^\ell}=\frac{1}{\beta_{2\ell}\cdot 6^\ell}
    =\frac{1}{\prod_{i=0}^{2\ell}\beta_i}.
\]
If $\beta_\ell=3$ then $y_{2\ell+1}\le 1$ by definition of the morphism $h_3$. Similarly, the above quantity is less than
\[
    \frac{1}{6^{\ell+1}}+\frac{1}{6^{\ell+1}}
    =\frac{2}{2\cdot 3 \cdot 6^\ell}
    =\frac{1}{\beta_{2\ell}\cdot 6^\ell}
    =\frac{1}{\prod_{i=0}^{2\ell}\beta_i}.
\]
\end{proof}

\end{document}
